% Referenzprojekt: https://www.overleaf.com/project/51fa85c3db89c3c351085071
% aus Tutorial-Serie: https://de.overleaf.com/learn/latex/How_to_Write_a_Thesis_in_LaTeX_(Part_1):_Basic_Structure

\documentclass[12pt]{article}
\usepackage[utf8]{inputenc}
% Deutsch: https://de.overleaf.com/learn/German
% "`	Left double quotes
% "'	Right double quotes
\usepackage[T1]{fontenc}
\usepackage[ngerman]{babel}
\usepackage{csquotes}

% documentation: https://ctan.kako-dev.de/macros/latex/contrib/biblatex/doc/biblatex.pdf
\usepackage[style=authoryear,sorting=nyt]{biblatex}
\addbibresource{references.bib}

% um Grafiken einzubinden, {Angabe des Pfades der Bilder}
\usepackage{graphicx}
\graphicspath{ {images/} }
% Nützliches Online Tool zur Tabellengenerierung: TablesGenerator.com 

% zur korreten Platzierung der Bilder im Titelblatt
\usepackage{chngpage}
\usepackage{calc}
\usepackage{geometry}

% Automatisches Generieren von Hyperlinks bei Verweisen (ref / cite)
\usepackage{hyperref}
\hypersetup{
  colorlinks   = true,    % Colours links instead of ugly boxes
  urlcolor     = blue,    % Colour for external hyperlinks
  linkcolor    = black,    % Colour of internal links
  citecolor    = red      % Colour of citations
}

% Infos über Dokument
\title{Ich habe starke schmerzen }
\author{Joe Mama}
\date{\today}

\begin{document}

    \newgeometry{top=2cm,bottom=2cm}

\begin{titlepage}
    
    \begin{flushright}
        \includegraphics[width=0.3\textwidth]{banner_top}
    \end{flushright}

    \begin{flushleft}
        \vspace*{1cm}
        \Huge
        \textbf{Homomorphe Verschlüsselung – Theorie und Praxis}
        
        \vspace{1.5cm}
        \textbf{Marvin Hahn}\\
        \textbf{Julian van Dyken}
        
        \vspace{0.5cm}
        \Large
        Exposé für die Lehrveranstaltung\\
        Spezielle Verfahren der IT-Sicherheit
        
        \vspace{0.5cm}
        im Studiengang\\
        Informatik
        
    \end{flushleft}        

    \vspace{0.5cm}
    \begin{flushleft}
        Hochschule Emden/Leer\\
        Prof. Dr. Patrick Fehlke\\
        M.Eng. Frederik Gosewehr\\
        \today
    \end{flushleft}    

\end{titlepage}

\restoregeometry
    \newpage

    \pagenumbering{Roman}   
    \tableofcontents
    \newpage

    \pagenumbering{arabic} 
    \section{Einleitung}
        \label{sec:einführung}
Normalerweise muss man verschlüsselte Daten entschlüsseln um damit Operationen 
durchzuführen z.B. Aufaddieren der Stimmen bei einer Wahl oder Daten für Maschinelles 
lernen. Der Nutzer muss daher einen Teil seiner Privatsphäre einbüßen, damit diese 
Verarbeitet werden können. Das birgt Risiken, weil die Informationen wer für wen 
gestimmt hat sehr sensibel sind und man auch nicht unbedingt Google seine sensiblen 
Daten für Maschinelles lernen zur Verfügung stellen möchte. Mit Homomorpher Verschlüsselung 
gibt es die Möglichkeit solche Operationen direkt auf den verschlüsselten Daten durchzuführen. 
Dadurch bleiben sensible Informationen geschützt. Auf die Beispiele von eben bezogen bedeutet 
das, dass es möglich ist ein Modell zu haben, was die Daten Homomorph auswertet, dass bedeutet, 
man gibt seine Informationen verschlüsselt an einen Anbieter, dieser kann diese dann auswerten 
und Berechnungen ausführen. Das Ergebnis dieser Berechnungen ist dann immer noch verschlüsselt. 
Dann bekommt man das Ergebnis zurück und kann es einfach entschlüsseln. Es ist also so, als ob 
die Berechnungen direkt auf dem Entschlüsselten Daten durchgeführt wurden.
    
    \section{Das Ziel der Ausarbeitung} 
        \subsection{Ziel}
    Grundlegende Ausarbeitung der Theorie (grundlegende Idee) hinter Homomorpher Verschlüsselung und deren Praktische Anwendung. Dabei gehen wir insbesondere auf die verschiedenen Arten der Homomorphen Verschlüsselung ein (partial, somewhat, full). Dabei werden wir beantworten, was genau diese Arten bedeutet und Beispiele für praktische Anwendungen geben. Dazu wird es eine Bewertung, ob und welche dieser Arten tatsächlich Praktisch einsetzbar sind geben, auf Basis der aktuellen Forschung.

    
    \section{Centralized vs Distributed}
        Versionskontrollsysteme werden hauptsächlich in zentrale (von engl. \textit{centralized}) und verteilte (von engl. \textit{distributed}) Systeme unterschieden. Daneben gibt es noch 
lokale Versionskontrolle\footcite[Siehe][S. 1]{pro_git}, die sich nur auf einen Rechner und meist nur ein Dokument beschränkt.

\subsection{Centralized}
	Zentralte Versionskontrollsysteme sind mit dem Client-Server abgebildet, wobei der Server meist nur Schnappschüsse, sprich die 
	Versionen von Dateien speichert\footcite[Vgl.][S. 671]{art:transition} und die Beteiligten sich nur mit dem Server synchronisieren.

\subsection{Distributed}
	Der Begriff \textit{verteilt} bedeutet in diesem Kontext, dass alle Beteiligten durch lokale Repositories die komplette Historie aller Dateien haben und sich untereinander synchronisieren (Peer-to-Peer).\\
	Man spricht nicht mehr von Versionen, sondern von Revisions, also Änderungen an einer Datei (zb. neue Zeile 10: "Hallo Welt").
	\footcite[Siehe][für anschauliche Beispiele]{online:intro}

\subsection{Vor-/Nachteile}
	Hier werden die jeweiligen Vor- und Nachteile der beiden Systeme betrachtet und als Übersicht in Trade-offs verpackt.\\
	Beispiel zentrale Systeme: Abhängigkeit eines Servers dafür mehr Kontrolle des Workflows (zb. durch Locking\footcite[Erklärung des Mechanisumses][S. 3-4]{book:svn})\\
	Beispiel verteilte Systeme: Es läuft offline und schnell dafür keine "`neuste Version"'\footcite[Vgl.][]{online:intro}

    
    \section{Praxis}
        \label{sec:praxis}
\begin{itemize}
    \item \textbf{Welche Theorien gibt es?}
    \begin{itemize}
        \item Grundlegende Theorie: partial, somewhat, full
        \item Bootstrapping
    \end{itemize}

    \item \textbf{Wo wollen wir anknüpfen?}
    \begin{itemize}
        \item Zunächst am Beispiel bekannter Methoden (z.B. RSA oder ElGamal) die Theorie erklären
    \end{itemize}
\end{itemize}
    
    \section{Fazit}
        Hier wird die Arbeit zusammengefasst und wieder auf die Stärken der beiden Arten eingegangen um ihre jeweilige Daseinsberechtigung zu untermauern.
    
    \newpage
    \printbibliography

\end{document}